\section{Head Integrated Gradients for Attention Head Attribution}
\label{ch:method}

Building on the Integrated Gradients background in Section~\ref{sec:preliminaries}, we now develop a head-attribution formulation for diffusion language models (DLLMs) and other mask-predict sequence models. Because training and evaluation in these models are defined by masked reconstruction over a supervised completion span rather than autoregressive next-token prediction, the attribution target must be aligned with the masked reconstruction objective.

To make head-level attribution amenable to integrated gradients, we introduce a differentiable gate for each selected attention head and define head importance through the resulting gated reconstruction objective. Our basic framework is called \emph{Head Integrated Gradients} (HeadIG). In the final method, we aggregate attributions over multiple paths in gate space; we refer to this multi-path variant as \emph{PolyHeadIG}, where ``Poly'' highlights the use of multiple sampled paths rather than a single synchronized trajectory.

Consider one example with token sequence $\mathbf{x} = \mathbf{c} \mathbin{\|} \mathbf{a} \in \mathcal{V}^{L}$, where $\mathbf{c}\in\mathcal{V}^{L_c}$ is the context prefix, $\mathbf{a}\in\mathcal{V}^{L_a}$ is the completion suffix, $\mathbin{\|}$ denotes sequence concatenation, and $L=L_c+L_a$. We write $x_i$ for the token at position $i$ in $\mathbf{x}$ and $\mathcal{C}=\{L_c+1,\dots,L\}$ for the completion span. Let $z(\mathbf{x})\in\mathbb{R}^{L\times |\mathcal{V}|}$ denote the model logits and $p_{i,v}(\mathbf{x})=\mathrm{softmax}(z_{i,:}(\mathbf{x}))_v$ the predicted probability of token $v\in\mathcal{V}$ at position $i$.

DLLMs supervise reconstruction only at masked positions. To better match DLLM denoising dynamics and reduce variance, we average over multiple masking rates $\{\rho_t\}_{t=1}^{T}$ and $S$ mask samples for each rate. For the $s$-th sampled mask $\mathbf{m}_{t,s}\in\{0,1\}^{L_a}$ at rate $\rho_t$, where $m_{t,s,j}=1$ indicates that the $j$-th completion token is masked and supervised, define the supervised index set $\mathcal{M}_{t,s}=\{\,L_c+j \mid j\in\{1,\dots,L_a\},\ m_{t,s,j}=1\,\}$ and the corresponding masked input $\tilde{\mathbf{x}}_{t,s}$ by
\begin{equation}
    \tilde{x}_{t,s,i}=
    \begin{cases}
        \texttt{[MASK]}, & i\in\mathcal{M}_{t,s},\\
        x_i, & \text{otherwise.}
    \end{cases}
\end{equation}
We require $|\mathcal{M}_{t,s}|\ge 1$ for every masking variant to avoid degenerate zero-loss samples.

We now instantiate attribution in attention-head space by inserting a differentiable gate at the output of each selected attention head. For a fixed masking variant $(t,s)$, let $\mathbf{O}^\ell\in\mathbb{R}^{L\times H\times d_h}$ denote the stacked per-head attention outputs at layer $\ell$ before the output projection, so that $O^\ell_{i,h,:}$ is the output vector of head $h$ at position $i$. We introduce a gate vector $\mathbf{g}_\ell\in\mathbb{R}^{H}$ and define the gated head output
\begin{equation}
    \tilde{O}_{i,h,:}^\ell(\mathbf{g}_\ell)
    =
    g_{\ell,h}\,O_{i,h,:}^\ell.
    \label{eq:gate_def}
\end{equation}
The remainder of the forward pass, including head concatenation, output projection, residual connections, and subsequent layers, remains unchanged. With a mild abuse of notation, we write $p_{i,v}(\tilde{\mathbf{x}}_{t,s};\mathbf{g})$ for the token probability induced by the gated forward pass on masking variant $(t,s)$. The corresponding masked reconstruction loss is
\begin{equation}
    \mathcal{L}_{t,s}(\mathbf{g})
    =
    \frac{1}{|\mathcal{M}_{t,s}|}
    \sum_{i\in\mathcal{M}_{t,s}} -\log p_{i,x_i}(\tilde{\mathbf{x}}_{t,s};\mathbf{g}).
\end{equation}
All model parameters are frozen; gradients are taken only with respect to the gate variables.

In practice, we gate all heads jointly over a layer range $\ell\in\{\ell_{\min},\dots,\ell_{\max}\}$ and flatten the corresponding gates into a single vector $\mathbf{g}\in\mathbb{R}^{D}$, where $D=(\ell_{\max}-\ell_{\min}+1)H$. Averaging the gated loss over masking rates and mask samples gives
\begin{equation}
    \mathcal{L}(\mathbf{g})
    =
    \frac{1}{TS}\sum_{t=1}^{T}\sum_{s=1}^{S}\mathcal{L}_{t,s}(\mathbf{g}).
\end{equation}
We define the attribution objective as $J(\mathbf{g})=-\mathcal{L}(\mathbf{g})$, so that larger values correspond to better reconstruction. Relative to the standard IG setup in Section~\ref{sec:preliminaries}, the input coordinates are now head-gate variables rather than raw input features. We use the all-zero gate as the baseline and the all-one gate as the target, i.e., $\boldsymbol{\gamma}(0)=\mathbf{0}_{D}$ and $\boldsymbol{\gamma}(1)=\mathbf{1}_{D}$. Motivated by the limitation of synchronized diagonal IG, we consider any differentiable path $\boldsymbol{\gamma}:[0,1]\to[0,1]^D$ connecting these two endpoints in gate space.

\paragraph{Definition (HeadIG Along a Path).}
For any such path $\boldsymbol{\gamma}$ and coordinate $d\in\{1,\dots,D\}$, we define the pathwise HeadIG attribution by
\begin{equation}
    \mathrm{IG}_{d}[\boldsymbol{\gamma}]
    =
    \int_{0}^{1}
    \frac{\partial J(\boldsymbol{\gamma}(\tau))}{\partial g_d}
    \frac{d\gamma_d(\tau)}{d\tau}\, d\tau.
    \label{eq:ig_general}
\end{equation}
This formulation preserves the IG principle while moving attribution from input space to head-gate space. It also makes clear that, once different heads follow different schedules along the path, the discrete estimator must weight each gradient by the corresponding coordinate-wise gate increment. The next section specifies how the path in gate space is chosen and how multiple sampled paths are averaged in our final method.